\chapter{Trabalhos Relacionados}

A integração do BIM com tecnologias de inteligência artificial generativa (IAG), como o ChatGPT, tem sido amplamente explorada no setor de arquitetura, engenharia e construção (AEC). Este tema abrange aplicações, frameworks inovadores, desafios e possibilidades futuras, sendo objeto de estudo de diversos trabalhos recentes que discutiremos nesta seção.

\begin{itemize}    
    \item \textbf {Integrating Building Information Modelling (BIM) with ChatGPT, Bard, and similar generative artificial intelligence in the architecture, engineering, and construction industry: applications, a novel framework, challenges, and future scope}: 
    
    \begin{itemize}
        \item   Um dos estudos analisados aborda a integração do BIM com modelos de IAG, como o ChatGPT e o Bard, no setor AEC. Este artigo identifica que a utilização dessas tecnologias visa otimizar processos como design e tomada de decisão. No entanto, os autores destacam desafios como falta de padronização de dados, deficiências em segurança e a necessidade de capacitação de profissionais do setor. Também apontam que os modelos generativos enfrentam dificuldades em interpretar a semântica complexa de dados BIM, além de problemas gerais como alta demanda computacional e falta de diretrizes éticas \cite{rane2023}.
        \item   Neste contexto, os pesquisadores propõem um framework para superar tais lacunas, com foco na otimização de processos de design, redução de custos e tempo, além de facilitar a adoção de IA por meio de interfaces intuitivas \cite{rane2023}.
    \end{itemize}
    
    \item \textbf {Prototyping a Chatbot for Site Managers Using Building Information Modeling (BIM) and Natural Language Understanding (NLU) Techniques}: 
    \begin{itemize}
        \item   Esse estudo relevante aborda o desenvolvimento de um chatbot, denominado JULE, projetado para gerentes de obras. Este modelo utiliza BIM e técnicas de processamento de linguagem natural (PLN) para fornecer informações precisas e em tempo real no canteiro de obras. O chatbot permite que os usuários consultem dados como dimensões de componentes estruturais e ferramentas necessárias de forma ágil e intuitiva, eliminando a dependência de interfaces convencionais baseadas em cliques e digitação \cite{lin2023}.
        \item   A abordagem proposta demonstrou alta precisão na identificação de intenções e na recuperação de informações essenciais, melhorando a adoção de sistemas baseados em BIM e promovendo maior eficiência no gerenciamento de obras \cite{lin2023}.
    \end{itemize}

    \item \textbf {Responsible AI in Construction Safety: Systematic Evaluation of Large Language Models and Prompt Engineering}: 
    \begin{itemize}
        \item   Esse trabalho investigou a aplicação de LLMs, como GPT-3.5 e GPT-4, em projetos BIM focados em segurança. Os autores realizaram testes com 385 questões de múltipla escolha de exames certificados pela BCSP, abrangendo áreas como gestão de sistemas de segurança, identificação de perigos e ergonomia. Utilizando técnicas de engenharia de \textit{prompt}, como \textit{Prompt Direct}, \textit{Prompt Chain-of-Thought} e \textit{Prompt Few-Shot}, os modelos foram avaliados em termos de precisão, consistência de respostas e erros categorizados. 
        \item   Os resultados indicaram que o GPT-4 alcançou 84,6\% de precisão, enquanto o GPT-3.5 obteve 73,8\%. Apesar do bom desempenho em áreas como gerenciamento de sistemas e controle de perigos, ambos os modelos apresentaram dificuldades em respostas a emergências e cálculos matemáticos. A pesquisa destaca o potencial desses modelos, mas também a necessidade de melhorias para evitar informações imprecisas ou enviesadas \cite{sammou2024}.
    \end{itemize}

    \item \textbf {Optimizing the Utilization of Generative Artificial Intelligence (AI) in the AEC Industry: ChatGPT Prompt Engineering and Design}:
    \begin{itemize}
        \item   Esse estudo explorou diferentes técnicas de engenharia de prompt para otimizar a utilização de LLMs no setor AEC. Os autores identificaram que métodos genéricos de interação limitam a eficácia desses modelos, resultando em saídas desalinhadas com as necessidades do setor. Problemas como confiabilidade, adaptação a regulamentações e aceitação no setor também foram identificados como barreiras \cite{samsami2024}. 
        \item   A pesquisa coletou técnicas de prompt eficazes, como prompts instrucionais e baseados em restrições, e testou essas abordagens em projetos reais. Os resultados mostraram que prompts específicos reduziram erros em cronogramas e melhoraram o reconhecimento de perigos. Apesar disso, limitações como alucinações do modelo para tarefas complexas e a necessidade de validação externa ainda persistem \cite{samsami2024}.
    \end{itemize}
\end{itemize}

Os trabalhos analisados fornecem uma visão abrangente sobre o estado da arte na integração de LLMs e BIM no setor AEC. Embora avanços significativos tenham sido alcançados, desafios importantes, como a interpretação semântica de dados BIM, segurança e a adaptação de LLMs às demandas específicas do setor, ainda necessitam de soluções mais robustas. Assim, este campo continua a ser promissor para futuras pesquisas e inovações.